\section{Metodologia}
Este trabalho implementou uma abordagem experimental para avaliação comparativa dos modelos preditivos descritos
anteriormente. Inicialmente, os dados do SNIRH foram submetidos a um pré-processamento abrangente, incluindo tratamento de
valores faltantes mediante técnicas de interpolação temporal, detecção de \textit{outliers} via análise por PCA e ECOD, e
normalização. Os dados foram particionados cronologicamente em conjuntos de treino (70\%), validação (15\%) e
teste (15\%), assegurando-se a preservação da sequência temporal para prevenir vazamento de informação.

Cada categoria de modelo exigiu transformações específicas dos dados, visto que cada um conta com um algoritmo
diferente de implementação e uma lógica interna distinta. Para o SARIMA, foi aplicada diferenciação e decomposição
sazonal. Modelos baseados em árvores requeriram codificação de variáveis temporais. Já redes LSTM demandaram janelamento
sequencial. O treinamento foi conduzido em \textit{pipelines} isolados, incorporando validação cruzada temporal e
otimização Bayesiana de hiperparâmetros para minimizar o KGE no conjunto de validação.

Na fase de avaliação, as predições geradas para o período de teste foram confrontadas com dados observados reservados de
teste mediante métricas hidrologicamente relevantes - coeficiente de Nash-Sutcliffe (NSE), RMSE e erro absoluto médio
(MAE) - complementadas por análise de sensibilidade a eventos extremos e avaliação de robustez computacional. Esta
estrutura metodológica visou garantir comparação objetiva do desempenho preditivo.

% Os resultados da análise comparativa serão disponibilizados através de um painel interativo (dashboard) desenvolvido
% após a consolidação das avaliações dos modelos. Esta plataforma de visualização, hospedada em acesso aberto, integrará
% dados históricos do nível do Rio Guaíba com as predições geradas pelos diferentes modelos, apresentando múltiplos
% cenários hidrológicos contextualizados por análises explicativas. Os usuários poderão explorar dinamicamente os dados
% mediante filtros temporais personalizáveis, permitindo focalizar períodos de interesse específico como eventos extremos
% ou fases sazonais críticas.

\section{Fonte de Dados}
Este trabalho empregou dados primários do Portal Hidroweb, integrante do Sistema Nacional de Informações sobre Recursos
Hídricos (SNIRH) da Agência Nacional de Águas (ANA) \cite{snirh}. Este sistema consolida informações de estações
hidrometeorológicas, incluindo níveis fluviais, vazões, precipitação, entre outros. A coleta automatizada dos dados foi
viabilizada mediante acesso à API do Hidroweb, concedido pela ANA ao autor.

As possíveis estações foram obtidas do endpoint \textit{HidroInventarioEstacoes} \cite{HidroInventarioEstacoes},
enquanto as séries detalhadas foram obtidas do endpoint \textit{HidroinfoanaSerieTelemetricaDetalhada}
\cite{HidroinfoanaSerieTelemetricaDetalhada}.

As estações analisadas pertencem à Bacia Hidrográfica do Guaíba, no Rio Grande do Sul, abrangendo em especial o Rio
Guaíba ('87200000') e os principais rios afluentes — Taquari ('86001000'), Caí ('87220000'), Gravataí ('87240000'),
Sinos ('87230000') e Jacuí ('85001000')— a fim de garantir ampla cobertura espacial e boa qualidade dos dados. Para
isto, foi necessário investigar todas as possíveis estações para estes rios. O total inicial era de 258 estações. Após
filtrar por estações em operação e estações telemétricas, restaram apenas 43 estações. As séries históricas destas foram
analisadas quanto à disponibilidade, qualidade, lacunas, duração e representatividade. Também excluíram-se estações com
finalidades específicas (energéticas, piezométricas etc.), que não apresentam as informações relevantes para este
estudo. 

\begin{figure}[htb] \caption{\label{funil} Funil de estações selecionadas do SNIRH} \begin{center}
  \includegraphics[scale=0.34]{img/funil_estações.png} \end{center} \legend{Fonte: Reprodução própria} \end{figure} 

Com isto, determinou-se que somente 15 estações na Bacia Hidrográfica do Guaíba apresentavam informações da forma
desejada. Destas, foram selecionadas pelo menos uma para cada um dos principais rios listados acima. Alguns rios, como o
Sinos e o Taquari, apresentavam 3 estações cujas informações poderiam ser utilizadas, mas por fim, foi determinado que
no máximo 2 estações seriam selecionadas de cada rio, a fim de manter um equilíbrio de dados e não introduzir um maior
peso para estes corpos hídricos, inclusive quando alguns, como o Rio Jacuí, só teriam 1 estação disponível.

Estas foram coletadas com dados à cada 15 minutos para um período histórico de 5 anos. Para o próprio Guaíba, adotou-se
como alvo, a estação CAIS MAUÁ C6 ('87450004') e complementarmente a estação do USINA DO GASÔMETRO - CHEIA 2024
('87444000'), estratégia necessária devido à descontinuidade nos registros da primeira estação em 2-3 de maio de 2024.
Conforme atestado pela SEMA \cite{informativo_sema}, este hiato decorreu de danos nos equipamentos de medição durante as
enchentes.

Por fim, a lista de estações selecionadas foram:
\begin{table}[h!]
    \centering
    \caption{Estações Hidrometeorológicas por Bacia}
    \label{tab:estacoes_hidro}
    \begin{tabular}{|l|l|l|}
        \hline
        \textbf{Rio} & \textbf{Estação} & \textbf{Código Estação} \\
        \hline
        \hline
        \multirow{2}{*}{\textbf{Guaíba}} & CAIS MAUÁ C6 & 87450004 \\
        \cline{2-3}
        & USINA DO GASÔMETRO (CHEIA 2024) & 87444000 \\
        \hline
        \multirow{2}{*}{\textbf{Taquari}} & MUÇUM & 86510000 \\
        \cline{2-3}
        & ENCANTADO & 86720000 \\
        \hline
        \multirow{2}{*}{\textbf{Caí}} & LINHA GONZAGA & 87150000 \\
        \cline{2-3}
        & BARCA DO CAÍ & 87170000 \\
        \hline
        \textbf{Gravataí} & PASSO DAS CANOAS - AUXILIAR & 87399000 \\
        \hline
        \multirow{2}{*}{\textbf{Sinos}} & SÃO LEOPOLDO & 87382000 \\
        \cline{2-3}
        & CAMPO BOM & 87380000 \\
        \hline
        \textbf{Jacuí} & RIO PARDO & 85900000 \\
        \hline
    \end{tabular}
\end{table}

Para uniformizar as séries, adotou-se `2018-08-01` como data inicial comum e `2025-12-01` como data final comum,
totalizando cerca de seis anos de dados telemétricos, totalizando mais de 2 milhões de registros. Devido ao limite de 30
dias imposto pela API, o processo de coleta foi realizado em chunks mensais posteriormente concatenados.

Foram coletados todos os dados disponíveis para estas estações neste período, incluindo dados do nível da água (e.g.,
\textit{Cota\_Adotada}), dados meteorológicos (e.g., \textit{Chuva\_Adotada}), dados descritivos da estação (e.g.,
\textit{Altitude}) e indicadores de controle de qualidade de cada variável, identificados pelo sufixo \textit{\_Status}.
Estas variáveis de status acompanham as variáveis de medição e seguem um padrão determinado pelo SNIRH que foi expandido
para utilizar estas variáveis como indicadores não só da qualidade dos dados, mas também como indicadores de variáveis
faltantes ou preenchidas. Este padrão é dado de tal forma: Status 0=Normal, 1=Suspeito, 2=Ruim, 3=Muito ruim,
4=Preenchido/Faltante, 5=Valor Atípico (outliers);

% Bateria, Chuva_Acumulada, Chuva_Acumulada_Status, Chuva_Adotada, Chuva_Adotada_Status, Cota_Adotada,
% Cota_Adotada_Status, Cota_Display, Cota_Display_Status, Cota_Manual, Cota_Manual_Status, Cota_Sensor,
% Cota_Sensor_Status, Data_Atualizacao, Data_Hora_Medicao, Pressao_Atmosferica, Pressao_Atmosferica_Status,
% Temperatura_Agua, Temperatura_Agua_Status, Temperatura_Interna, Vazao_Adotada, Vazao_Adotada_Status, codigoestacao,
% Temperatura_Interna_Status

\section{Processamento de dados}
O primeiro e maior desafio enfrentado nos dados coletados foram os valores faltantes. Certos atributos estão
consistentemente faltando do dataset, como valores de 'Pressao\_Atmosferica' e 'Temperatura\_Agua' infelizmente
apresentam muitos valores faltantes para serem de qualquer uso nessa análise e logo serão desconsiderados daqui pra
frente.

\begin{figure}[htb] \caption{\label{faltantes} Valores faltantes} \begin{center}
  \includegraphics[scale=0.53]{img/faltantes.png} \end{center} \legend{Fonte: Reprodução própria} \end{figure} 

Também vale considerar todos aqueles relacionados à Cota. Estes, em especial, são formas de medir o nível da água.
'Cota\_Manual', por exemplo, representa medidas feitas com uma régua ao invés das medidas automatizadas do sensor. Dado
isto, usaremos valores de 'Cota\_Manual' e 'Cota\_Sensor' para preencher valores faltantes de 'Cota\_Adotada';

Para escolher quais dados seriam utilizados nos modelos de predição, também foi realizada uma análise da correlação dos
valores, presente na \autoref{correlacao}.

\begin{figure}[htb] \caption{\label{correlacao} Correlação entre os dados das estações} \begin{center}
  \includegraphics[scale=0.275]{img/correlacao.png} \end{center} \legend{Fonte: Reprodução própria} \end{figure} 

A análise de correlação evidencia que as variáveis 'Cota\_Manual', 'Cota\_Sensor' e 'Cota\_Display' apresentam forte
correlação linear com o alvo 'Cota\_Adotada'. Observa-se que 'Cota\_Manual' possui mais de 99\% de valores ausentes, e
os registros válidos restantes concentram-se majoritariamente em pontos nos quais 'Cota\_Adotada' não está disponível.
Esse padrão sugere que as medições manuais são utilizadas como mecanismo de contingência em situações de manutenção ou
falha dos sensores automáticos. Adicionalmente, 'Cota\_Sensor' e 'Cota\_Display' mostram-se, na maior parte do tempo,
virtualmente idênticas a 'Cota\_Adotada', o que é consistente com o fluxo operacional da estação, no qual as leituras
oficiais derivam predominantemente dos sensores automáticos.

A variável 'Vazao\_Adotada' apresenta elevada correlação com 'Cota\_Adotada', o que é consistente com o comportamento
hidrodinâmico esperado: incrementos de vazão (m\textsuperscript{3}/s) em um canal de geometria aproximadamente fixa
tendem a elevar o nível d'água até o ponto em que o escoamento interage áreas de várzea adjacentes. Essa variável
constitui um preditor importante devido à sua forte associação com a variável alvo e à sua baixa correlação com o
restante dos atributos, indicando aporte de informação complementar ao modelo. No entanto, há uma alta taxa de valores
ausentes (>15\%), cuja magnitude varia entre estações. A estação alvo (Guaíba) não possui registros de vazão, o que
inviabiliza seu uso direto nessa localidade. Ainda assim, dada a elevada dependência física entre vazão e nível,
optou-se por manter essa variável no processo de modelagem, mesmo que as mesmas impliquem um maior processamento dos
dados.

Com base na análise de correlação e na avaliação de valores faltantes, foram selecionadas as seguintes variáveis para os
modelos de \textit{Machine Learning}: 'Cota\_Adotada', correspondente ao nível do rio em centímetros, cujos valores
ausentes serão imputados a partir de 'Cota\_Manual' e 'Cota\_Sensor' sempre que disponíveis; 'Chuva\_Adotada',
representando a precipitação acumulada desde a última leitura, em milímetros; 'Temperatura\_Interna', Correspondente à
temperatura atmosférica medida no equipamento, em °C; e 'Vazao\_Adotada', que indica o fluxo de água registrado pelo
sensor em $\frac{m^{3}}{s}$.

A variável 'Chuva\_Acumulada' também foi inicialmente escolhida e trabalhada junto das outras variáveis por um bom
tempo. Contudo, depois de análises mais aprofundadas, a parte "Acumulada" da mesma se mostrou inconsistênte, dado que
deveria representar a precipitação acumulada na estação por mês, mas o mesmo não se mostrou consistênte através das
diferentes estações. Logo, junto de outras transformações, a mesma foi recriada a partir da variável 'Chuva\_Adotada',
em um processo que será explicado melhor mais a frente.

\section{Ferramentas e recursos utilizados}
O desenvolvimento da parte escrita deste trabalho será feita em \LaTeX, utilizando o modelo \abnTeX, desenvolvido por
Araujo et al (2015) \cite{abntex2classe} de acordo com as regras da ABNT NBR 14724:2011, ABNT NBR 6024:2012 e outras.

A implementação dos modelos será desenvolvida em Python, selecionado por seu ecossistema robusto para análise de dados e
aprendizado de máquina. A biblioteca \textit{statsmodels} será empregada para modelagem SARIMA, enquanto redes LSTM
serão construídas utilizando a API de baixo nível do \textit{TensorFlow}. Para os modelos baseados em árvores, serão
adotadas as bibliotecas nativas do XGBoost e LightGBM. Uma arquitetura modular será implementada mediante classes
especializadas para cada modelo, integrando \textit{pipelines} do \textit{scikit-learn} para garantir fluxos
reprodutíveis de pré-processamento. O processamento numérico e manipulação de séries temporais serão realizados com
\textit{NumPy} e \textit{Pandas}, respectivamente, assegurando eficiência computacional e tratamento dos dados.

O painel interativo utilizará \textit{Streamlit} para conversão ágil de scripts Python em aplicações web visuais,
garantindo integração nativa com bibliotecas científicas (\textit{Pandas}, \textit{Matplotlib}, \textit{Plotly}). A
hospedagem no \textit{Streamlit Community Cloud} assegurará disponibilidade pública contínua e atualizações automáticas
conforme novos dados ou modelos forem incorporados.

\section{Setup Experimental para treinamento de modelos} 
The training process was conducted in a machine with the following specifica- tions: 12th Gen Intel(R) Core(TM)
i5-12500H, 16GB RAM and a NVIDIA GeForce RTX 3050 Laptop 4GB VRAM. The LSTM model takes in training around 1.5 seconds
per epoch, which led to around 2.5 to 4 minutes per experiment. Additionally, the feature importance calculation for
each prediction takes around 1 second per prediction.